% Options for packages loaded elsewhere
\PassOptionsToPackage{unicode}{hyperref}
\PassOptionsToPackage{hyphens}{url}
\documentclass[
]{article}
\usepackage{xcolor}
\usepackage[margin=1in]{geometry}
\usepackage{amsmath,amssymb}
\setcounter{secnumdepth}{-\maxdimen} % remove section numbering
\usepackage{iftex}
\ifPDFTeX
  \usepackage[T1]{fontenc}
  \usepackage[utf8]{inputenc}
  \usepackage{textcomp} % provide euro and other symbols
\else % if luatex or xetex
  \usepackage{unicode-math} % this also loads fontspec
  \defaultfontfeatures{Scale=MatchLowercase}
  \defaultfontfeatures[\rmfamily]{Ligatures=TeX,Scale=1}
\fi
\usepackage{lmodern}
\ifPDFTeX\else
  % xetex/luatex font selection
\fi
% Use upquote if available, for straight quotes in verbatim environments
\IfFileExists{upquote.sty}{\usepackage{upquote}}{}
\IfFileExists{microtype.sty}{% use microtype if available
  \usepackage[]{microtype}
  \UseMicrotypeSet[protrusion]{basicmath} % disable protrusion for tt fonts
}{}
\makeatletter
\@ifundefined{KOMAClassName}{% if non-KOMA class
  \IfFileExists{parskip.sty}{%
    \usepackage{parskip}
  }{% else
    \setlength{\parindent}{0pt}
    \setlength{\parskip}{6pt plus 2pt minus 1pt}}
}{% if KOMA class
  \KOMAoptions{parskip=half}}
\makeatother
\usepackage{color}
\usepackage{fancyvrb}
\newcommand{\VerbBar}{|}
\newcommand{\VERB}{\Verb[commandchars=\\\{\}]}
\DefineVerbatimEnvironment{Highlighting}{Verbatim}{commandchars=\\\{\}}
% Add ',fontsize=\small' for more characters per line
\usepackage{framed}
\definecolor{shadecolor}{RGB}{248,248,248}
\newenvironment{Shaded}{\begin{snugshade}}{\end{snugshade}}
\newcommand{\AlertTok}[1]{\textcolor[rgb]{0.94,0.16,0.16}{#1}}
\newcommand{\AnnotationTok}[1]{\textcolor[rgb]{0.56,0.35,0.01}{\textbf{\textit{#1}}}}
\newcommand{\AttributeTok}[1]{\textcolor[rgb]{0.13,0.29,0.53}{#1}}
\newcommand{\BaseNTok}[1]{\textcolor[rgb]{0.00,0.00,0.81}{#1}}
\newcommand{\BuiltInTok}[1]{#1}
\newcommand{\CharTok}[1]{\textcolor[rgb]{0.31,0.60,0.02}{#1}}
\newcommand{\CommentTok}[1]{\textcolor[rgb]{0.56,0.35,0.01}{\textit{#1}}}
\newcommand{\CommentVarTok}[1]{\textcolor[rgb]{0.56,0.35,0.01}{\textbf{\textit{#1}}}}
\newcommand{\ConstantTok}[1]{\textcolor[rgb]{0.56,0.35,0.01}{#1}}
\newcommand{\ControlFlowTok}[1]{\textcolor[rgb]{0.13,0.29,0.53}{\textbf{#1}}}
\newcommand{\DataTypeTok}[1]{\textcolor[rgb]{0.13,0.29,0.53}{#1}}
\newcommand{\DecValTok}[1]{\textcolor[rgb]{0.00,0.00,0.81}{#1}}
\newcommand{\DocumentationTok}[1]{\textcolor[rgb]{0.56,0.35,0.01}{\textbf{\textit{#1}}}}
\newcommand{\ErrorTok}[1]{\textcolor[rgb]{0.64,0.00,0.00}{\textbf{#1}}}
\newcommand{\ExtensionTok}[1]{#1}
\newcommand{\FloatTok}[1]{\textcolor[rgb]{0.00,0.00,0.81}{#1}}
\newcommand{\FunctionTok}[1]{\textcolor[rgb]{0.13,0.29,0.53}{\textbf{#1}}}
\newcommand{\ImportTok}[1]{#1}
\newcommand{\InformationTok}[1]{\textcolor[rgb]{0.56,0.35,0.01}{\textbf{\textit{#1}}}}
\newcommand{\KeywordTok}[1]{\textcolor[rgb]{0.13,0.29,0.53}{\textbf{#1}}}
\newcommand{\NormalTok}[1]{#1}
\newcommand{\OperatorTok}[1]{\textcolor[rgb]{0.81,0.36,0.00}{\textbf{#1}}}
\newcommand{\OtherTok}[1]{\textcolor[rgb]{0.56,0.35,0.01}{#1}}
\newcommand{\PreprocessorTok}[1]{\textcolor[rgb]{0.56,0.35,0.01}{\textit{#1}}}
\newcommand{\RegionMarkerTok}[1]{#1}
\newcommand{\SpecialCharTok}[1]{\textcolor[rgb]{0.81,0.36,0.00}{\textbf{#1}}}
\newcommand{\SpecialStringTok}[1]{\textcolor[rgb]{0.31,0.60,0.02}{#1}}
\newcommand{\StringTok}[1]{\textcolor[rgb]{0.31,0.60,0.02}{#1}}
\newcommand{\VariableTok}[1]{\textcolor[rgb]{0.00,0.00,0.00}{#1}}
\newcommand{\VerbatimStringTok}[1]{\textcolor[rgb]{0.31,0.60,0.02}{#1}}
\newcommand{\WarningTok}[1]{\textcolor[rgb]{0.56,0.35,0.01}{\textbf{\textit{#1}}}}
\usepackage{graphicx}
\makeatletter
\newsavebox\pandoc@box
\newcommand*\pandocbounded[1]{% scales image to fit in text height/width
  \sbox\pandoc@box{#1}%
  \Gscale@div\@tempa{\textheight}{\dimexpr\ht\pandoc@box+\dp\pandoc@box\relax}%
  \Gscale@div\@tempb{\linewidth}{\wd\pandoc@box}%
  \ifdim\@tempb\p@<\@tempa\p@\let\@tempa\@tempb\fi% select the smaller of both
  \ifdim\@tempa\p@<\p@\scalebox{\@tempa}{\usebox\pandoc@box}%
  \else\usebox{\pandoc@box}%
  \fi%
}
% Set default figure placement to htbp
\def\fps@figure{htbp}
\makeatother
\setlength{\emergencystretch}{3em} % prevent overfull lines
\providecommand{\tightlist}{%
  \setlength{\itemsep}{0pt}\setlength{\parskip}{0pt}}
\usepackage{bookmark}
\IfFileExists{xurl.sty}{\usepackage{xurl}}{} % add URL line breaks if available
\urlstyle{same}
\hypersetup{
  pdftitle={Network meta-analysis},
  hidelinks,
  pdfcreator={LaTeX via pandoc}}

\title{Network meta-analysis}
\author{}
\date{\vspace{-2.5em}}

\begin{document}
\maketitle

\textbf{Inference lab:} Hypothesis → Visualise → Assumptions → Conduct →
Conclude.

\subsection{Learning objectives}\label{learning-objectives}

\begin{itemize}
\tightlist
\item
  Fit a network meta-analysis (NMA) using \texttt{netmeta}.
\item
  Draw and read the network graph.
\item
  Derive SUCRA rankings and interpret them cautiously.
\end{itemize}

\subsection{Prerequisites}\label{prerequisites}

Pairwise meta-analysis (W4 S2).

\subsection{Background}\label{background}

Network meta-analysis extends pairwise meta-analysis to more than two
treatments by combining direct comparisons (from head-to-head trials)
with indirect comparisons (inferred through a shared comparator). The
central assumption is \emph{transitivity}: the studies compared
indirectly are exchangeable --- no moderator distinguishes them in a way
that systematically shifts the effect. The statistical analogue,
\emph{consistency}, checks whether direct and indirect estimates agree.

SUCRA (surface under the cumulative ranking curve) summarises each
treatment's rank distribution on a 0-1 scale, where 1 is best. SUCRA
values compress a lot of uncertainty into a single number; they are
useful for signalling but should be reported with interval estimates for
every pairwise comparison.

\subsection{Setup}\label{setup}

\begin{Shaded}
\begin{Highlighting}[]
\FunctionTok{library}\NormalTok{(tidyverse)}
\FunctionTok{library}\NormalTok{(netmeta)}
\FunctionTok{set.seed}\NormalTok{(}\DecValTok{42}\NormalTok{)}
\FunctionTok{theme\_set}\NormalTok{(}\FunctionTok{theme\_minimal}\NormalTok{(}\AttributeTok{base\_size =} \DecValTok{12}\NormalTok{))}
\end{Highlighting}
\end{Shaded}

\subsection{1. Hypothesis}\label{hypothesis}

H₀: all treatments in the network have equal efficacy. H₁: at least two
differ.

\subsection{2. Visualise}\label{visualise}

\begin{Shaded}
\begin{Highlighting}[]
\FunctionTok{head}\NormalTok{(Senn2013)}

\NormalTok{net }\OtherTok{\textless{}{-}} \FunctionTok{netmeta}\NormalTok{(TE, seTE, treat1, treat2, studlab,}
               \AttributeTok{data =}\NormalTok{ Senn2013, }\AttributeTok{sm =} \StringTok{"MD"}\NormalTok{,}
               \AttributeTok{common =} \ConstantTok{FALSE}\NormalTok{, }\AttributeTok{random =} \ConstantTok{TRUE}\NormalTok{,}
               \AttributeTok{reference.group =} \StringTok{"plac"}\NormalTok{)}
\end{Highlighting}
\end{Shaded}

\begin{Shaded}
\begin{Highlighting}[]
\FunctionTok{netgraph}\NormalTok{(net, }\AttributeTok{plastic =} \ConstantTok{FALSE}\NormalTok{, }\AttributeTok{thickness =} \StringTok{"number.of.studies"}\NormalTok{,}
         \AttributeTok{points =} \ConstantTok{TRUE}\NormalTok{, }\AttributeTok{cex.points =} \DecValTok{2}\NormalTok{, }\AttributeTok{cex =} \FloatTok{0.8}\NormalTok{)}
\end{Highlighting}
\end{Shaded}

\subsection{3. Assumptions}\label{assumptions}

\begin{Shaded}
\begin{Highlighting}[]

\end{Highlighting}
\end{Shaded}

The within-design heterogeneity and between-design inconsistency
variances are the key diagnostics.

\subsection{4. Conduct}\label{conduct}

\begin{Shaded}
\begin{Highlighting}[]
\NormalTok{rank\_df }\OtherTok{\textless{}{-}} \FunctionTok{netrank}\NormalTok{(net, }\AttributeTok{small.values =} \StringTok{"good"}\NormalTok{)}
\NormalTok{rank\_df}
\end{Highlighting}
\end{Shaded}

\subsection{5. Concluding statement}\label{concluding-statement}

\begin{Shaded}
\begin{Highlighting}[]
\NormalTok{ranked }\OtherTok{\textless{}{-}} \FunctionTok{as\_tibble}\NormalTok{(rank\_df}\SpecialCharTok{$}\NormalTok{ranking.random, }\AttributeTok{rownames =} \StringTok{"treatment"}\NormalTok{) }\SpecialCharTok{|\textgreater{}}
  \FunctionTok{rename}\NormalTok{(}\AttributeTok{sucra =}\NormalTok{ value) }\SpecialCharTok{|\textgreater{}}
  \FunctionTok{arrange}\NormalTok{(}\FunctionTok{desc}\NormalTok{(sucra))}
\NormalTok{top }\OtherTok{\textless{}{-}} \FunctionTok{head}\NormalTok{(ranked, }\DecValTok{1}\NormalTok{)}
\end{Highlighting}
\end{Shaded}

\begin{quote}
In a network meta-analysis of
\texttt{length(unique(c(Senn2013\$treat1,\ Senn2013\$treat2)))}
treatments across \texttt{length(unique(Senn2013\$studlab))} studies
(\texttt{nrow(Senn2013)} pairwise comparisons), the treatment with the
highest SUCRA for HbA1c lowering was \textbf{\texttt{top\$treatment}}
(SUCRA = \texttt{round(top\$sucra,\ 2)}). SUCRA rankings should be
interpreted alongside the pairwise effect estimates and their intervals,
not as a standalone conclusion.
\end{quote}

\subsection{Common pitfalls}\label{common-pitfalls}

\begin{itemize}
\tightlist
\item
  Treating SUCRA as if it were the effect size.
\item
  Ignoring the transitivity assumption because the output compiled.
\item
  Mixing direct and indirect evidence without quantifying inconsistency.
\end{itemize}

\subsection{Further reading}\label{further-reading}

\begin{itemize}
\tightlist
\item
  Rücker G, Schwarzer G. (2015). \emph{Ranking treatments in frequentist
  network meta-analysis.}
\item
  Salanti G. (2012). \emph{Indirect and mixed-treatment
  comparison\ldots{}}
\end{itemize}

\subsection{Session info}\label{session-info}

\begin{Shaded}
\begin{Highlighting}[]

\end{Highlighting}
\end{Shaded}

\subsection{See also --- chapter index}\label{see-also-chapter-index}

\begin{itemize}
\tightlist
\item
  \href{../index.Rmd}{Methods in this chapter}
\item
  \href{../../../ch00-find-your-method.Rmd}{Find your method (Chapter
  0)}
\end{itemize}

\end{document}
